In conclusion, our project was able to build a heat production improvement system with the capability to automatically schedule production units for utility company HeatItOn. The system allowed for low operational costs while producing heat and satisfying customer demand which resulted in a rise in overall profits. It considered electricity price, seasonal heat demand fluctuations, maintenance downtimes and production unit operating capacity limits. The team's final solution consisted of an intuitive frontend visualization, intelligent backend optimization logic, and structured data storage, all developed with a flexible architecture. By utilizing Scrum methods through the project, the team was able to iterate development, receive frequent stakeholder feedback, and gradually improve team technical abilities as well as organizational processes. Also by doing tests, reviewing code and making improvements the team improved the quality of the code, made it easier to maintain, increased system's reliability and enhanced teamwork during the development process.\par

Eventually the team's solution addressed the problem by changing a manual scheduling process into an automated optimization algorithm. The optimizer schedules when to produce heat to meet the heating demand while keeping the operational costs low and maximizing profits in the electricity market. The application lets users easily see detailed optimization data, configure production units to their liking, export data and check detailed optimization data. Also one of the needs for storing CSV results was changed while developing the project. The team found that using JSON for storage was more flexible and easier to handle for what we wanted to achieve. This allows the system to reduce human effort in creating schedules, minimize error, and efficiently allocate resources. Overall, this project proved that an automatic optimization tool can be helpful for planning operations in district heating systems.\par

Even though main goals have been achieved, there are still some parts of the project that were not completely finished or are still somewhat limited. The current version of the system relies on local JSON and CSV files rather than relational databases or distributed systems. Some technical debt stayed in certain areas of the code because of limitations in the framework and fast development cycles. Using Avalonia and LiveCharts for the frontend brought up some issues with keeping the code easy to manage, making the user interface responsive, and ensuring good performance when displaying graphics. Moreover, testing methods were brought in quite late during the development process, which led to uneven test coverage throughout the system. There were still some problems related to Scrum that had not been fixed. These problems included not judging how long tasks would take accurately and issues with team members not showing up and slowdowns in the work process during several sprints.\par

Future efforts could aim on improving both technical setup and usability of the app. A future version of the system might move to a more advanced setup that uses relational databases, cloud services and distributed systems. This would help it handle bigger datasets and support many users at the same time. The optimizer can be improved by adding more different parameters to take for optimization, better ways to manage environmental goals like CO2 emissions and the ability to run on multiple cores for enhanced performance. Extra features like user login, authorization, support for multiple languages, real-time heat map display, better reporting tools and improved scenario management could make the system even more useful. Future development should focus more on automated testing, continuous integration and building a frontend architecture that is easy to maintain. This will help make the system more reliable and make sure it can grow over time.\par
